\section{Threats to Validity}
\label{sec:threats}
This section catalogs threats that can invalidate the Paper~B2 investigation even when Protocol v1.0 is followed. Threats are stated conservatively and are bounded to the evaluated cohort and the admitted evidence representations.

\subsection{Internal validity}
\begin{itemize}
  \item \textbf{Extraction and transcription errors:} BOR entries may contain mistakes when transcribing facts from sources.
  \item \textbf{Pipeline drift:} If any extraction/normalization tooling is used, changes in tool versions or parameters can alter DER/MSR outputs.
  \item \textbf{Trace breakage:} Missing or ambiguous trace pointers can allow unsupported conclusions.
\end{itemize}

\subsection{Construct validity}
\begin{itemize}
  \item \textbf{Operationalization mismatch:} The protocol's operational definition of ``Reasoning,'' ``Legitimacy,'' and ``Execution'' may not align with each system's native terminology.
  \item \textbf{Proxy measurement error:} Measurements may be imperfect proxies for boundary separation.
  \item \textbf{Schema pressure:} Normalizing heterogeneous systems into shared DER object types can obscure meaningful differences.
\end{itemize}

\subsection{External validity}
\begin{itemize}
  \item \textbf{Cohort limitations:} Findings are restricted to the selected governance/authorization cohort and its admitted representations.
  \item \textbf{Temporal validity:} Systems evolve; observations tied to specific versions or documentation snapshots may not generalize to future releases.
  \item \textbf{Deployment variance:} Paper~B2 analyzes system representations, not a randomized sample of real deployments.
\end{itemize}

\subsection{Selection bias}
The cohort is not randomly sampled. Inclusion may over-represent widely documented or widely deployed systems and under-represent niche or proprietary governance stacks.

\textbf{Selection rationale.} The cohort is purposive and coverage-oriented: it spans cloud IAM, policy-as-code, Kubernetes admission/RBAC, service-mesh authorization, relationship-based authorization, and secret/access governance systems represented by admitted public documentation.

\subsection{Observer bias}
Analysts may preferentially record observations that fit pre-existing intuitions about layered architectures.

\textbf{Observer-bias mitigation.} BOR, SRF, DER, MSR, dataset, analysis, retained-classification, and cohort-conclusion artifacts are separated and machine-validated; lack of a formal double-coding process remains an explicit limitation rather than a hidden claim of inter-analyst agreement.

\subsection{Researcher expectancy}
Because the candidate boundary is pre-specified, analysts may unintentionally steer derivations toward confirming structures.

\textbf{Mitigation stance.} Paper~B2 reduces expectancy risk by: (i) separating BOR (facts) from DER (derivations) and MSR (measurements), and (ii) requiring trace pointers for each conclusion.

\subsection{Reproducibility limitations}
\begin{itemize}
  \item \textbf{Source stability:} Vendor documentation may change; stable snapshots and hashes are required.
  \item \textbf{Tooling determinism:} If automated extraction/measurement is used, it must be pinned (lockfiles) and rerunnable.
  \item \textbf{Inter-analyst agreement:} Some derivations may lack formal coder-agreement support without a dedicated reconciliation plan.
\end{itemize}

\textbf{Replication package.} The replication surface is the repository itself: canonical JSON artifacts, schemas, build scripts, registry checks, publication manifest generation, and validation commands recorded for deterministic rerun.
